\setauthor{\secondauthor}
\subsection{Erfüllung der Zielsetzung}

Ziel der Diplomarbeit war es, zwei verschiedene Informationsstrategien zur 
digitalen Vermarktung der App \textit{Jobbing} zu entwickeln, technisch umzusetzen und 
an realen Nutzungsdaten miteinander zu vergleichen. Dieses Ziel konnte im wesentlichen erreicht werden.

Der Multi-Step-Funnel sowie der Chatbot konnten erfolgreich umgesetzt und 
testbar gemacht werden. Für beide Konzepte konnten auch Nutzungsdaten erhoben 
und ausgewertet werden, so dass nicht nur die technische Machbarkeit der beiden 
Ansätze überprüft, sondern auch ihre Wirkung auf Nutzer:innen analysiert werden konnte.

Ein weiterer Teil des Zieles war es, den Chatbot in mehreren Varianten 
umzusetzen, um den Einfluss der Personalisierung und der Tonalität zu 
untersuchen. Auch dieses Ziel wurde erreicht. Die erhobenen Daten zeigen, dass 
Personalisierung im Chatbot die wahrgenommene Relevanz positiv beeinflussen kann.
Die Wirkung der Tonalität ist weniger eindeutig und muss immer in bezug auf 
Gesprächsinhalt, Zielgruppe und generelle Nutzerführung betrachtet werden.

Darüber hinaus konnte auch die Frage beantwortet werden, welche der beiden Informationsstrategien im konkreten Projektkontext wirksamer ist. Hier zeigte die Auswertung, dass der Multi-Step-Funnel im Gesamtvergleich bessere Werte lieferte. Insbesondere die längere Verweildauer und das größere Download-Interesse lassen vermuten, dass die schlüssige und schrittweise Nutzerführung in dieser speziellen Situation die bessere Strategie war.

Auch im Bereich Datenerhebung und -auswertung konnte ein wesentliches Ziel der Arbeit erreicht werden. Es zeigte sich, dass weder reine Chatbot-Daten noch klassische Webanalyse für sich allein ausreichen, um das Nutzerverhalten umfassend zu erfassen. Erst die Kombination aus Voiceflow, Microsoft Clarity und Google Analytics ermöglichte eine fundierte und vergleichbare Auswertung beider Konzepte.

Insgesamt kann festgehalten werden, dass die Zielsetzung der Diplomarbeit erreicht wurde. Die Arbeit brachte technische und analytische Ergebnisse und damit eine datengetragene Beurteilung der beiden Informationsstrategien.

\subsection{Fehler}

Trotz der insgesamt positiven Ergebnisse traten im Zuge des
Projektes auch einige Einschränkungen und Schwächen auf, sowohl bei der
Umsetzung als auch bei der Aussagekraft der späteren Auswertung.
Ein wesentlicher Punkt ist die Datenlage. Die Anzahl der
ausgewerteten Sitzungen war zur Erkennung erster Tendenzen ausreichend,
jedoch nicht, um allgemeingültige Aussagen mit hoher statistischer
Sicherheit zu treffen. Die Ergebnisse sind also eher als fundierte
Projektresultate und weniger als übertragbare Marktanalyse zu
deuten.

Beim Chatbot zeigte sich zudem, dass die Wahrnehmung der Qualität der
Nutzererfahrung stark von der jeweils konkreten Formulierung der
Antworten abhängt. Teilweise wurden Antworten als zu
allgemein, zu unpersönlich oder als zu wenig konkret empfunden. Dies
veranschaulicht, dass ein Chatbot nur dann gut überzeugt, wenn nicht
nur die technische Struktur funktioniert, sondern auch die sprachliche
Ausgestaltung sehr sorgfältig ausjustiert ist. Besonders im Marketingkontext
kann schon eine leicht unpassende Tonalität die Wirkung des
Systems stark mindern.

Trotz dieser vergleichbaren Informationsstrategien war die Gegenüberstellung der beiden Ansätze nicht in allen Punkten ganz symmetrisch.
Der Chatbot lieferte uns wesentlich mehr qualitative Einblicke in die
konkreten Fragen und Unsicherheiten der Nutzer:innen, während der
Multi-Step-Funnel primär über klassische Webmetriken auszuwerten
war. Das führt dazu, dass sich die erhobenen Datenarten an manchen Stellen unterschiedlich sind
 und ein direkter Vergleich an einigen Stellen erschwert ist.

Ein weiterer Fehler bzw. Schwachpunkt des Projektes war, dass einige
Kennzahlen einfach zu isoliert betrachtet werden konnten. Werte wie
Sitzungsdauer, Nachrichtenanzahl oder Seitenaufrufe sind nur dann
sinnvoll interpretierbar, wenn sie in Verbindung mit Zielhandlungen,
Feedback und Gesprächsverlauf ausgewertet werden. Das wurde erst im
Laufe der Arbeit klar.

Insgesamt wird deutlich, dass die Ergebnisse sicherlich aussagekräftig sind,
jedoch immer im Rahmen der Projektbedingungen interpretiert werden
müssen. Vor allem die geringe Datenmenge und die hohe
Fehleranfälligkeit des Chatbots im Hinblick auf sprachliche Feinheiten
zeigen wesentliche Einschränkungen der Arbeit.

\subsection{Erweiterungs- und Verbesserungspotenzial}

Für beide Informationsstrategien ergeben sich verschiedene Anknüpfungspunkte 
zur Weiterentwicklung. Besonders beim Chatbot besteht noch Spielraum bei der 
inhaltlichen Treffsicherheit und der Qualität der gelieferten Antworten. Eine 
erweiterte Wissensbasis, besser abgestimmte Formulierungen und eine noch 
stärkere Führung durch den Dialog könnten dazu beitragen, die Zufriedenheit 
der Nutzer:innen und das Download-Interesse weiter zu steigern.

Zudem wäre es lohnenswert, die Personalisierung noch weiter 
zielgruppenorientiert auszubauen. Die bisherigen Ergebnisse zeigen bereits, 
dass personalisierte Ansprache positiv wirkt, in den nächsten Iterationen könnte 
dies differenzierter angegangen werden: mit spezifischeren Argumentationslinien 
für Gastronom:innen und Jobsuchende oder einer noch besseren Anpassung an 
individuelle Interessen innerhalb des Dialogs.

Auch beim Multi-Step-Funnel gibt es noch weitere Möglichkeiten zur Optimierung. 
Die Arbeit hat gezeigt, dass eine klare, strukturierte Nutzerführung einen positiven Einfluss hat. 
Daraus ergibt sich die Möglichkeit, in zukünftigen Iterationen noch gezielter 
einzelne Seiten, Informationsblöcke oder UI-Elemente zu testen. A/B-Tests zu 
unterschiedlichen Reihenfolgen, Formulierungen oder Call-to-Action-Elementen 
könnten zusätzliche Verbesserungen ermöglichen.

Ein wichtiger nächster Schritt wäre auch die Ausweitung der Testbasis. Je mehr Nutzer:innen man hat und je länger man testet, desto belastbarer werden die Ergebnisse. Dadurch ließen sich nicht nur die Wirkungen der einzelnen Chatbot-Varianten besser absichern, sondern auch die Überlegenheit des Multi-Step-Funnels.

Schließlich wäre auch denkbar, die Konzepte künftig nicht nur getrennt, sondern stärker ineinander verwoben zu betrachten. Denkbar wäre dabei eine Nutzung des Multi-Step-Funnels als primäre Informationsstrategie, während ein Chatbot ergänzend zum Beispiel bei Rückfragen, für individuelle Erklärungen oder die Klärung konkreter Unsicherheiten eingesetzt werden könnte. Beides zusammen könnte die Stärken beider Ansätze nutzen: die Übersichtlichkeit des Funnels und die Interaktivität des Chatbots.

Insgesamt zeigt die Arbeit, dass sowohl der Chatbot als auch der Multi-Step-Funnel über das aktuelle Projekt hinaus weiterentwickelt werden können. Dabei ist eine besonders wertvolle Erkenntnis, dass datenbasierte Evaluation nicht nur der abschließenden Bewertung dient, sondern auch die zentrale Grundlage für zukünftige Optimierungsschritte darstellt.
